%=======================================================================
%  CHAPTER 8 — WIRELESS SYSTEMS AND STANDARDS  (6 hrs)
%=======================================================================
\section{Wireless Systems and Standards}

{\color{sub}\footnotesize 6 hrs \gap$\cdot$\gap asked in \mk{21 of 22 papers}
(only 72 Ma skips it) \gap$\cdot$\gap 5 syllabus sub-topics \gap$\cdot$\gap
\mk{no numerical component}: the one calculation the syllabus touches, GSM frame
efficiency, is worked in \textbf{Ch 7 companion \S3}.}

\lead{Read this first:}
\begin{itemize}
  \item This chapter is \mk{always the last question on the paper}, and it is nearly always
        a full 8 marks plus a 4 or 5 mark rider. Two of its topics are among the six most
        repeated in the whole subject.
  \item Five blocks carry the chapter, in order of return:
  \begin{itemize}
    \item \textbf{Regulatory issues} \tS{9} $\rightarrow$ pure recall, no diagram, no
          formula. \mk{The cheapest 4 to 5 marks in the subject.}
    \item \textbf{GSM traffic and control channels} \tS{8} $\rightarrow$ the BCH / CCCH /
          DCCH tree, with what each channel does and which way it points.
    \item \textbf{GSM system architecture} \tS{7} $\rightarrow$ one diagram plus one line
          per block.
    \item \textbf{IS-95 forward and reverse channels} \tS{6} $\rightarrow$ draw the
          composite block diagram and name the four forward, two reverse channels.
    \item \textbf{GSM frame structure and hierarchy} \tF{5} $\rightarrow$ hyperframe down
          to time slot, with the four durations.
  \end{itemize}
  \item \mk{Learn 4 diagrams cold}: GSM architecture, GSM frame hierarchy, the forward
        CDMA (IS-95) channel, and the GSM signal-processing chain. They answer about
        30 marks.
  \item \mk{Name clash to watch.} In this chapter \textbf{BCH $=$ broadcast channel} and
        \textbf{SCH $=$ synchronisation channel}. In \textbf{Ch 6} BCH is the
        Bose--Chaudhuri--Hocquenghem block code. Same letters, unrelated.
  \item GSM's access scheme (TDMA/FDMA/FDD with slow hopping) is argued in \textbf{Ch 7};
        the frame efficiency proof is in \textbf{Ch 7 companion \S3}. \mk{Do not re-derive
        either here}, cross-refer and spend the words on the channels.
  \item \textbf{WLAN and WiFi} are answered in \textbf{Ch 7 \S7.11}. \mk{WiMAX, LTE, UMTS
        and EVDO belong here}, in \S8.13.
\end{itemize}

\hr

\T{8.1 GSM: what it is, its services and its features}

\Q{What is GSM? State its services and features. Give its specifications}{\m{5}}
{\color{sub}\footnotesize \yr{\textbf{74 Bh} \m{5} (specifications of GSM)}
\gap$\cdot$\gap never asked alone as a definition, but \mk{every GSM question opens with
these three lines}, so they are worth two minutes}

\sbsr{0.52}{%
\lead{Definition and history}
\begin{itemize}
  \item \textbf{GSM} $=$ global system for mobile communications.
  \item \mk{The first cellular system to specify digital modulation, network-level
        architecture and services together}, not just an air interface.
  \item \textbf{1982}: CEPT sets up \emph{Groupe Sp\'ecial Mobile} to build one
        pan-European cellular system in the 900 MHz band.
  \item \textbf{1990}: the same specification is applied at 1800 MHz as \textbf{DCS 1800},
        a PCN deployment started in the UK.
  \item \textbf{1992}: renamed \textbf{Global System for Mobile Communications}, for
        marketing reasons. The acronym stayed.
  \item Later carried into the USA at 1900 MHz as \textbf{PCS 1900}.
\end{itemize}}{%
\lead{Features worth naming \m{2}}
\begin{itemize}
  \item \textbf{Superior speech quality} against the analog systems it replaced.
  \item \textbf{High security}: SIM-based authentication and over-the-air ciphering.
  \item \textbf{International roaming}, because one standard covers many countries.
  \item \textbf{Low-power hand-held terminals}, helped by discontinuous transmission.
  \item \textbf{Open, published standard} $\rightarrow$ many vendors, low equipment cost.
\end{itemize}}

\vspace{2pt}
\sbsr{0.52}{%
\lead{Three categories of service}
\begin{enumerate}
  \item \textbf{Teleservices}: voice, emergency calling, \textbf{SMS}, fax.
  \item \textbf{Bearer services}: data transport, \mk{300 to 9600 bps}.
  \item \textbf{Supplementary services}: call forwarding, call barring, call waiting,
        caller ID, conferencing.
\end{enumerate}}{%
\lead{The three frequency plans}
\par\vspace{2pt}
\begin{tabularx}{\linewidth}{@{}L{1.9cm}XX@{}}
\toprule
\hrow \thead{System} & \thead{Uplink (MS $\rightarrow$ BS)} & \thead{Downlink (BS $\rightarrow$ MS)} \\
\midrule
GSM 900  & 890--915 MHz   & 935--960 MHz \\
DCS 1800 & 1710--1785 MHz & 1805--1880 MHz \\
PCS 1900 & 1850--1910 MHz & 1930--1990 MHz \\
\bottomrule
\end{tabularx}
\vspace{2pt}
{\color{sub}\footnotesize All three are \textbf{FDD}. \mk{Uplink is always the lower band},
because path loss is slightly lower there and the handset is the power-limited end.
\yr{SST Ch 8}}}

\T{8.2 GSM system architecture \; \small\mk{one of the two most repeated questions in this chapter}}

\Q{Draw and explain the GSM architecture; explain the operation of each component}{\m{4}\m{5}\m{7}\m{8}\m{4+4}}
{\color{sub}\footnotesize \tS{7} \yr{\textbf{\texttt{82 Bh}} \m{4} (NSS only),
\texttt{80 Ba} \m{8}, \textbf{\texttt{79 Bh}} \m{8}, \textbf{73 Bh} \m{4} (NSS) $+$ \m{5},
73 Ma \m{7}, 71 Ma \m{4+4}, 70 Ma \m{8}} \gap$\cdot$\gap
\mk{four papers ask only for the NSS}, so learn that box in more detail than the rest}

\sbsr{0.50}{%
\figT{c8_arch.png}
\vspace{1pt}
{\color{sub}\footnotesize \textbf{Draw exactly this.} Three dashed boxes, three named
interfaces, and the public networks hanging off the MSC. \yr{SST Ch 8}}

\vspace{3pt}
\lead{1. Base Station Subsystem (BSS), the radio subsystem}
\begin{itemize}
  \item \textbf{MS} $=$ mobile station $=$ \textbf{ME} (mobile equipment) $+$ \textbf{SIM}.
  \begin{itemize}
    \item ME carries the \textbf{IMEI}, the equipment's own identity.
    \item SIM carries the \textbf{IMSI}, the subscriber's identity, plus the secret key.
    \item \mk{Identity is in the card, not the handset} $\rightarrow$ personal mobility.
  \end{itemize}
  \item \textbf{BTS} $=$ base transceiver station. Houses the radio transceivers that
        define one cell. Talks to the MS over the \textbf{Um} air interface.
  \item \textbf{BSC} $=$ base station controller.
  \begin{itemize}
    \item Manages radio resources for one to several hundred BTSs.
    \item Radio channel setup, frequency hopping, power control.
    \item \mk{Performs handover between two BTSs under itself}; anything wider goes to
          the MSC.
    \item Connects to the BTS over the \textbf{Abis} interface.
  \end{itemize}
\end{itemize}}{%
\lead{2. Network and Switching Subsystem (NSS) \m{4}}
\begin{itemize}
  \item Handles \textbf{switching of GSM calls between the BSCs and external networks},
        and holds every subscriber database.
  \item \textbf{MSC} $=$ mobile switching centre. The central unit: controls traffic among
        BSCs, sets up and releases calls, handles inter-BSC handover, and gateways to
        \textbf{PSTN, ISDN and data networks}. Reaches the BSC over the \textbf{A}
        interface.
  \item \textbf{HLR} $=$ home location register.
  \begin{itemize}
    \item Permanent record of every subscriber \textbf{homed} on this MSC: IMSI, number,
          subscribed services.
    \item Also holds the subscriber's \textbf{current location}, so an incoming call can be
          routed to wherever the mobile now is.
  \end{itemize}
  \item \textbf{VLR} $=$ visitor location register.
  \begin{itemize}
    \item Temporarily stores the IMSI and customer data of every \textbf{roaming}
          subscriber inside this MSC's area.
    \item On log-in the MSC tells the visitor's \textbf{home} HLR where they are, so the
          home HLR can route incoming calls over the PSTN.
  \end{itemize}
  \item \textbf{AuC} $=$ authentication centre. Strongly protected database of the
        authentication and ciphering keys of every subscriber in the HLR and VLR.
  \item \textbf{EIR} $=$ equipment identity register.
  \begin{itemize}
    \item Lists IMEIs as white (allowed), grey (tracked) or black (barred).
    \item A stolen handset's IMEI (dial \code{*\#06\#} to read it) is blacklisted here, and
          optionally in the \textbf{CEIR}, which bars it on every operator that subscribes.
    \item \mk{EIR bars the handset, AuC authenticates the subscriber.} Different targets.
  \end{itemize}
\end{itemize}

\lead{3. Operation Support Subsystem (OSS) \m{2}}
\begin{itemize}
  \item Built around one or more \textbf{OMC} $=$ operation and maintenance centres.
  \item Lets engineers \textbf{monitor, diagnose and troubleshoot} every part of the
        network: MS, BTS, BSC, MSC.
  \item Maintains hardware and network performance records.
  \item Handles \textbf{charging and billing}.
\end{itemize}}

\vspace{2pt}
\sbsr{0.47}{%
\figT{c8_interfaces.png}
\vspace{1pt}
{\color{sub}\footnotesize The three standardised interfaces. \mk{Naming them is usually
worth a mark on its own}, and they are what makes multi-vendor GSM possible.
\yr{SST Ch 8, Fig 11.6}}}{%
\lead{The interfaces, and why they matter}
\begin{itemize}
  \item \textbf{Um} $=$ MS to BTS, the \textbf{air interface}. The only radio hop.
  \item \textbf{Abis} $=$ BTS to BSC. Standardised, so a BTS and a BSC may come from
        different vendors.
  \item \textbf{A} $=$ BSC to MSC. Standardised.
  \item \textbf{SS7} carries the signalling on from the MSC into the PSTN.
  \item \mk{Because Abis and A are open, an operator can mix vendors.} That is the whole
        commercial point of GSM being a \emph{network} standard and not only an air
        interface.
\end{itemize}

\lead{Answer skeleton \m{8}}
\begin{enumerate}
  \item Draw the figure. Three subsystems boxed, interfaces labelled. \m{3}
  \item BSS: MS, BTS, BSC, one line each. \m{2}
  \item NSS: MSC, HLR, VLR, AuC, EIR, one line each. \m{2}
  \item OSS: OMC, monitoring and billing. \m{1}
\end{enumerate}}

\T{8.3 GSM radio subsystem and specifications}

\Q{Specifications of GSM; the radio subsystem and ARFCN}{\m{5}}
{\color{sub}\footnotesize \yr{\textbf{74 Bh} \m{5}} \gap$\cdot$\gap
\mk{these numbers also feed the frame-structure and frame-efficiency answers}, so they earn
their place even though the question itself is rare}

\sbsr{0.53}{%
\lead{The radio plan}
\begin{itemize}
  \item Two \textbf{25 MHz} bands: \textbf{890--915 MHz reverse} (MS to BS),
        \textbf{935--960 MHz forward} (BS to MS).
  \item Multiple access is \textbf{FDD $+$ TDMA over FDMA}, with optional slow
        \textbf{frequency hopping}. See \textbf{Ch 7 \S7.3}.
  \item Each band is cut into \textbf{200 kHz} carriers called \textbf{ARFCN}s, absolute
        radio frequency channel numbers.
  \item \mk{One ARFCN names a forward and reverse pair 45 MHz apart}, so an ARFCN is a
        duplex channel, not a single frequency.
  \item $25\ \text{MHz} / 200\ \text{kHz} = 125$ carriers, but a \textbf{100 kHz guard band}
        is left at each edge, so only \mk{124 duplex channels} are usable, ARFCN 1 to 124.
  \item Each 200 kHz carrier is time shared by \textbf{8 users} using TDMA.
  \item \textbf{Physical channel $=$ (ARFCN, time slot number)}. A \textbf{logical channel}
        is then mapped onto a physical channel, see \S8.5.
  \item DCS 1800 uses 1710--1785 / 1805--1880 MHz, three times the spectrum, giving
        \textbf{374} duplex channels.
\end{itemize}}{%
\figT{c8_gsm_specs.png}
\vspace{1pt}
{\color{sub}\footnotesize \textbf{Reproduce this table for a \m{5} specifications question.}
\yr{Rappaport Table 10.3, via SST Ch 8} \gap$\cdot$\gap
\mk{one printed value to know about}: the table gives the voice coder at \textbf{13.4 kbps},
while the RPE-LTP codec itself produces \textbf{260 bits per 20 ms $=$ 13.0 kbps} exactly
(\S8.7). Quote 13 kbps if you are asked for the codec rate. \added}

\vspace{3pt}
\lead{Rates worth memorising}
\begin{itemize}
  \item Modulation \textbf{0.3 GMSK}, channel data rate \mk{270.833 kbps} per carrier.
  \item Bandwidth efficiency $= 270.833 / 200 = \mathbf{1.354}$ bps/Hz.
  \item Per-user channel rate $= 270.833 / 8 = \mathbf{33.854}$ kbps.
  \item After all overhead, \mk{maximum user data rate 24.7 kbps}.
  \item Bit period \mk{3.692 $\mu$s}, time slot \mk{576.9 $\mu$s}, frame \mk{4.615 ms}.
\end{itemize}}

\T{8.4 GSM frame structure and frame hierarchy}

\Q{Draw and explain the GSM frame structure / channel structure / frame hierarchy}{\m{4}\m{5}\m{6}}
{\color{sub}\footnotesize \tF{5} \yr{\textbf{\texttt{82 Bh}} \m{4} (frame hierarchy),
\textbf{77 Ch} \m{4}, \textbf{76 Bh} \m{5}, 74 Ma \m{6} (channel structure),
\textbf{70 Bh} \m{4}} \gap$\cdot$\gap
74 Ma adds \mk{``show that TDMA frame efficiency cannot reach 100\,\% in GSM''} \m{2},
\textbf{worked in full in Ch 7 companion \S3}}

\sbsr{0.52}{%
\figT{c8_frame_hierarchy.png}
\vspace{1pt}
{\color{sub}\footnotesize \textbf{Draw the hierarchy top down and label the four durations.}
That is the whole answer to a \m{4} or \m{5} question. \yr{SST Ch 8}}}{%
\lead{The five levels, bottom up}
\begin{enumerate}
  \item \textbf{Time slot} $=$ \mk{156.25 bits $=$ 576.92 $\mu$s}. One burst.
  \item \textbf{TDMA frame} $=$ 8 slots $=$ 1250 bits $=$ \mk{4.615 ms} $\left(=120/26\
        \text{ms}\right)$.
  \item \textbf{Multiframe}, two kinds:
  \begin{itemize}
    \item \textbf{26-frame traffic multiframe} $=$ \mk{120 ms}. Carries TCH, SACCH, FACCH.
    \item \textbf{51-frame control multiframe} $=$ \mk{235.4 ms}. Carries BCH and CCCH.
  \end{itemize}
  \item \textbf{Superframe} $=$ 51 traffic multiframes $=$ 26 control multiframes
        $=$ 1326 TDMA frames $=$ \mk{6.12 s}.
  \item \textbf{Hyperframe} $=$ 2048 superframes $=$ 2\,715\,648 TDMA frames
        $=$ \mk{3 h 28 min 53.76 s}.
\end{enumerate}
\vspace{2pt}
{\color{sub}\footnotesize \mk{Why a hyperframe at all}: the ciphering algorithm and the
hopping sequence are both driven by the frame number, so the frame number must not repeat
for a long time. \gap$\cdot$\gap \mk{Two multiframe sizes, one superframe}: $51 \times 26$
either way $=$ 1326, which is exactly why 26 and 51 were chosen. \added
\gap$\cdot$\gap The SST slide prints the hyperframe as 3 h 28 min \emph{52.76} s;
$2048 \times 6.12 = 12533.76$ s recomputes to \mk{53.76 s}. \added}}

\vspace{2pt}
\sbsr{0.47}{%
\figT{c8_slot_breakdown.png}
\vspace{1pt}
{\color{sub}\footnotesize One frame down to one normal burst. The \mk{546.5 $\mu$s} is the
useful part, \mk{577 $\mu$s} is the whole slot including guard space. \yr{SST Ch 8}}}{%
\lead{Inside a normal burst \m{3}}
\begin{itemize}
  \item \textbf{3 tail bits} $(0,0,0)$ $\rightarrow$ mark the start, let the equalizer
        settle.
  \item \textbf{57 encrypted data bits}.
  \item \textbf{1 stealing flag} $\rightarrow$ says whether this half is user data or
        \textbf{FACCH} signalling.
  \item \textbf{26 training bits} $\rightarrow$ a known sequence the receiver correlates
        against to \mk{build a channel model for the equalizer} (Ch 5).
  \item \textbf{1 stealing flag}, \textbf{57 encrypted data bits}, \textbf{3 tail bits}.
  \item \textbf{8.25 guard bits} $\rightarrow$ stop bursts from overlapping and give the
        power amplifier ramp-up and ramp-down time.
  \item Total $3+57+1+26+1+57+3 = \mathbf{148}$ useful $+\ 8.25$ guard $=$
        \mk{156.25 bits}.
\end{itemize}}

\vspace{2pt}
\sbsr{0.55}{%
\figT{c8_burst_formats.png}
\vspace{1pt}
{\color{sub}\footnotesize \textbf{Five burst formats, all 156.25 bits.} \yr{Rappaport
Fig 10.9, via SST Ch 8}}}{%
\lead{Why five formats}
\begin{itemize}
  \item \textbf{Normal burst} $\rightarrow$ TCH and most control channels.
  \item \textbf{FCCH burst} $\rightarrow$ 142 fixed bits, \mk{all zeros}. Under GMSK an
        all-zero block becomes a pure tone, which is exactly what a mobile needs to pull its
        local oscillator onto the base station.
  \item \textbf{SCH burst} $\rightarrow$ a long \mk{64-bit training sequence}, because the
        mobile has no channel estimate yet.
  \item \textbf{RACH (access) burst} $\rightarrow$ 8 start bits and a \mk{68.25-bit extended
        guard period}. The mobile does not yet know its timing advance, so the burst is kept
        short to guarantee it cannot land on the next slot.
  \item \textbf{Dummy burst} $\rightarrow$ mixed bits, no information. Sent on unused slots
        of the BCH carrier so its \mk{power stays constant} and neighbouring mobiles can
        measure it.
\end{itemize}}

\vspace{2pt}
\sbs{%
\figT{c8_speech_multiframe.png}
\vspace{1pt}
{\color{sub}\footnotesize \textbf{Traffic (speech) multiframe, 26 frames, 120 ms.}
\yr{SST Ch 8, Fig 11.7}}

\vspace{2pt}
\lead{Traffic multiframe \m{2}}
\begin{itemize}
  \item TCH data runs in \textbf{consecutive frames}, but \mk{frames 13 and 26 are not TCH}.
  \item \textbf{Frame 13} carries \textbf{SACCH}.
  \item \textbf{Frame 26} carries \textbf{idle} bits for a full-rate TCH, and SACCH for a
        half-rate TCH.
  \item The idle frame is what gives the mobile time to \mk{measure neighbouring cells} for
        MAHO (Ch 2).
\end{itemize}}{%
\figT{c8_control_multiframe.png}
\vspace{1pt}
{\color{sub}\footnotesize \textbf{Control multiframe, 51 frames, 235 ms.} (a) forward link
on TS0, (b) reverse link on TS0. \yr{SST Ch 8, Fig 11.8}}

\vspace{2pt}
\lead{Control multiframe \m{2}}
\begin{itemize}
  \item \textbf{Forward}: F (FCCH) then S (SCH) repeat \mk{every 10 frames}; between them sit
        B (BCCH) and C (PCH / AGCH) bursts, with one idle frame.
  \item \textbf{Reverse}: all 51 frames are \textbf{R}, reverse RACH bursts. \mk{The uplink
        control channel is pure random access}, so no structure is possible.
  \item Both live in \textbf{TS0 of the BCH ARFCN}.
\end{itemize}}

\vspace{2pt}
\lead{Frame efficiency, in one line \m{2}}
\begin{itemize}
  \item $\eta_f = \left(1 - \dfrac{b_{OH}}{b_T}\right)\times 100\,\%$, with
        $b_T = 8 \times 156.25 = 1250$ bits and $b_{OH} = 8\,(6 + 8.25 + 26) = 322$ bits
        $\rightarrow \eta_f = \mathbf{74.24\,\%}$.
  \item \mk{It can never be 100\,\%}: tail, training and guard bits are needed for the
        equalizer and for timing, and removing them breaks the link.
  \item \textbf{Full working, with the 74 Ma and 70 Bh statements, is in Ch 7
        companion \S3.} Do not redo it here.
\end{itemize}

\T{8.5 GSM logical channels: traffic and control \; \small\mk{the most repeated question in this chapter}}

\Q{Explain the working of all traffic and control channels used in GSM}{\m{4}\m{5}\m{6}\m{8}\m{4+4}}
{\color{sub}\footnotesize \tS{8} \yr{\textbf{\texttt{81 Bh}} \m{8} (control channels),
\textbf{\texttt{80 Bh}} \m{4+4} (BCH $+$ CCCH), \textbf{79 Ch} \m{5} (BCCH only),
\textbf{78 Ch} \m{8}, \textbf{77 Ch} \m{6}, \textbf{75 Bh} \m{8}, \textbf{72 Ash} \m{8},
\textbf{70 Bh} \m{4}} \gap$\cdot$\gap
\mk{80 Bh asks for BCH and CCCH only}, 79 Ch for BCCH only. Read which one is wanted}

\sbsr{0.42}{%
\figT{c8_logical_channels.png}
\vspace{1pt}
{\color{sub}\footnotesize \textbf{Draw this tree first, then write one line per leaf.}
\yr{SST Ch 8}}

\vspace{3pt}
\lead{Physical against logical}
\begin{itemize}
  \item \textbf{Physical channel} $=$ one ARFCN $+$ one time slot number. Hardware.
  \item \textbf{Logical channel} $=$ a stream of a particular kind, \mk{mapped onto} a
        physical channel. Meaning.
  \item Many logical channels time-share the same physical channel over the multiframe.
\end{itemize}}{%
\lead{Traffic channels (TCH) \m{2}}
\begin{itemize}
  \item Carry \textbf{digitally encoded user speech or user data}, both directions.
  \item \textbf{Full rate TCH/FS}: a slot in \mk{every} frame $\rightarrow$ 22.8 kbps.
  \item \textbf{Half rate TCH/HS}: a slot in \mk{every second} frame $\rightarrow$
        11.4 kbps, doubling the users per carrier.
  \item Data variants: TCH/F9.6, F4.8, F2.4 and TCH/H4.8, H2.4.
\end{itemize}

\lead{1. Broadcast channels (BCH) \; forward only, on TS0 of the BCH ARFCN \m{4}}
\begin{itemize}
  \item \textbf{BCCH}, broadcast control channel.
  \begin{itemize}
    \item Broadcasts \textbf{cell and network identity}, the list of channels in use, and the
          access parameters.
    \item \mk{This is what a nearby mobile reads to identify a cell and lock on to it.}
  \end{itemize}
  \item \textbf{FCCH}, frequency correction channel.
  \begin{itemize}
    \item Occupies TS0 of the first frame, \mk{repeated every 10th frame} of the control
          multiframe.
    \item Lets each subscriber unit \textbf{synchronise its local oscillator} to the base
          station's exact frequency.
  \end{itemize}
  \item \textbf{SCH}, synchronisation channel.
  \begin{itemize}
    \item Broadcast in TS0 \textbf{immediately after an FCCH frame}.
    \item Identifies the serving base station and gives \textbf{frame synchronisation}; the
          base station also issues \textbf{timing advance} here.
  \end{itemize}
\end{itemize}}

\vspace{2pt}
\sbsr{0.46}{%
\lead{2. Common control channels (CCCH) \; call establishment \m{4}}
\begin{itemize}
  \item Occupy TS0 of every control frame \textbf{not} used by the BCH or the idle frame.
  \item \textbf{PCH}, paging channel, \emph{forward}: notifies a specific mobile of an
        \textbf{incoming call}.
  \item \textbf{RACH}, random access channel, \emph{reverse}: the mobile
        \textbf{acknowledges a page} or \textbf{originates a call}. \mk{Contended, so
        collisions are possible} and a random backoff is used.
  \item \textbf{AGCH}, access grant channel, \emph{forward}: the base station's answer to a
        RACH. Assigns a time slot, radio channel and dedicated control channel.
        \mk{AGCH is the last CCCH message before the mobile leaves the control channel.}
\end{itemize}}{%
\lead{3. Dedicated control channels (DCCH) \; both directions \m{4}}
\begin{itemize}
  \item Bidirectional, same format and function each way. May sit in \textbf{any slot on
        any ARFCN except TS0 of the BCH ARFCN}.
  \item \textbf{SDCCH}, stand-alone dedicated control channel.
  \begin{itemize}
    \item Signalling during idle periods and \textbf{during call setup, before a TCH is
          allocated}; released as soon as the TCH is assigned.
    \item Carries authentication, ciphering setup, location update and SMS.
    \item \mk{The intermediate channel that keeps mobile and base connected while resources
          are being allocated.}
  \end{itemize}
  \item \textbf{SACCH}, slow associated control channel.
  \begin{itemize}
    \item \mk{Always paired with a TCH or an SDCCH}; occupies \textbf{every 13th frame} of
          the traffic multiframe.
    \item \emph{Forward}: power-level and timing-advance commands.
    \item \emph{Reverse}: received signal strength and quality reports for the serving cell
          \textbf{and its neighbours} $\rightarrow$ this is what feeds MAHO.
  \end{itemize}
  \item \textbf{FACCH}, fast associated control channel.
  \begin{itemize}
    \item Carries \textbf{urgent} signalling, above all \textbf{handover} commands.
    \item Has no slots of its own: it \mk{steals a TCH burst}, flagged by the two stealing
          bits in the normal burst.
    \item \mk{Far faster than SACCH}, at the cost of 20 ms of speech, which the speech
          decoder conceals.
  \end{itemize}
\end{itemize}}

\T{8.6 Example of a GSM call, and location updating}

\Q{GSM frame structure; how traffic and control channels are used while making a call}{\m{4+4}\m{4+6}}
{\color{sub}\footnotesize \tF{2} \yr{\textbf{70 Bh} \m{4+4}, \textbf{77 Ch} \m{4+6}}
\gap$\cdot$\gap \mk{the frame-structure half is \S8.4}; this band answers the second half
\gap$\cdot$\gap the channel names must come out in the right \emph{order}, that is what is
being marked}

\sbsr{0.55}{%
\lead{Mobile-originated call setup, in order}
\begin{enumerate}
  \item MS reads \textbf{FCCH} $\rightarrow$ pulls its oscillator onto the base station.
  \item MS reads \textbf{SCH} $\rightarrow$ frame synchronisation and base station identity.
  \item MS reads \textbf{BCCH} $\rightarrow$ cell and network parameters. It is now
        \mk{camped on the correct BCH}.
  \item To originate, MS sends a burst of \textbf{RACH} data on the \textbf{same ARFCN} it
        is locked to.
  \item Base station answers with an \textbf{AGCH} message on the CCCH, assigning an
        \textbf{SDCCH}.
\end{enumerate}
\vspace{2pt}
{\color{sub}\footnotesize \mk{Mobile-terminated differs in one step}: the network first pages
the mobile on the \textbf{PCH}, the mobile answers on the \textbf{RACH}, and everything from
step 6 is identical. \added}}{%
\figT{c8_location_update.png}
\vspace{1pt}
{\color{sub}\footnotesize \textbf{Location updating, channel by channel}: which logical
channel carries each step when a mobile is switched on in a new area. Note how much of it
rides on the \textbf{SDCCH}. \yr{SST Ch 8, Table 2.1}}}

\vspace{2pt}
\sbs{%
\begin{enumerate}
  \setcounter{enumi}{5}
  \item MS, which is monitoring \textbf{slot 0 of the BCH}, reads its ARFCN and slot
        assignment from the AGCH and \textbf{retunes to the SDCCH}.
  \item Over the SDCCH the MS receives \textbf{timing advance} and \textbf{power level}
        commands. Only now can it send proper normal bursts.
  \item SDCCH carries \textbf{authentication} messages both ways. Meanwhile the PSTN
        connects the dialled party to the MSC, and the MSC switches the speech path to the
        serving base station.
  \item Base station commands the MS over the SDCCH to a \textbf{new ARFCN and slot for the
        TCH}.
  \item MS tunes to the \textbf{TCH}; speech flows both ways and the \mk{SDCCH is vacated}.
\end{enumerate}}{%
\lead{What each stage costs, and why it is built this way}
\begin{itemize}
  \item Steps 1--3 are \mk{listen only}. The mobile transmits nothing until it has read the
        BCH, which is why a handset can be switched on inside a foreign network safely.
  \item Step 4 is the \mk{only contended transmission} in the whole sequence, so it uses the
        short RACH burst with its extended guard period (\S8.4).
  \item Step 7 exists because the mobile does not yet know its distance from the base
        station. \mk{Without timing advance its normal burst would overlap the next slot.}
  \item The SDCCH does \mk{all the work of call setup} and is then thrown away. That is what
        ``stand-alone dedicated'' means: dedicated to one mobile, standing alone from any
        traffic channel.
  \item Once on the TCH, further signalling has to \mk{steal bursts (FACCH)} or wait for the
        13th frame (SACCH), because there is no spare capacity left.
\end{itemize}}

\T{8.7 Signal processing in GSM: speech in, radio out}

\Q{Basic signal processing operations to convert speech to a radio signal and back}{\m{8}}
{\color{sub}\footnotesize \tP{2} \yr{\textbf{71 Bh} \m{8}, \texttt{81 Ba} \m{8}}
\gap$\cdot$\gap \mk{always the full 8 marks} \gap$\cdot$\gap
speech coding itself is \textbf{Ch 6}, channel coding and interleaving are \textbf{Ch 6},
equalization is \textbf{Ch 5}. \mk{Here you assemble them into one chain}}

\sbsr{0.55}{%
\figT{c8_signal_proc.png}
\vspace{1pt}
{\color{sub}\footnotesize \textbf{Draw the forward chain, then say the reverse is its mirror.}
The radio channel sits between modulation and demodulation. \yr{Rappaport Fig 10.11, via
SST Ch 8}}

\vspace{3pt}
\lead{The chain, block by block}
\begin{enumerate}
  \item \textbf{Digitizing and source coding}: RPE-LTP speech coder,
        \mk{260 bits per 20 ms $=$ 13 kbps}.
  \item \textbf{Channel coding}: adds redundancy, \mk{260 $\rightarrow$ 456 bits}
        $=$ 22.8 kbps.
  \item \textbf{Interleaving}: the 456 bits are cut into \mk{eight 57-bit sub-blocks} and
        spread over \textbf{8 consecutive TCH slots}.
  \item \textbf{Burst formatting}: adds tail, training, stealing and guard bits
        $\rightarrow$ the 156.25-bit burst of \S8.4.
  \item \textbf{Ciphering}: A5 stream cipher, keystream XORed with the data bits.
  \item \textbf{Modulation}: \mk{0.3 GMSK at 270.833 kbps}.
  \item Radio channel, then \textbf{demodulation, de-ciphering, burst formatting,
        de-interleaving, channel decoding, source decoding} $\rightarrow$ speech.
\end{enumerate}}{%
\lead{Speech coding \m{2}}
\begin{itemize}
  \item Coder is a \textbf{RELP / RPE-LTP}: residually excited linear predictor with a long
        term predictor. Full detail in \textbf{Ch 6}.
  \item Output \mk{260 bits per 20 ms frame}. Half-rate codec runs at 6.5 kbps.
  \item \textbf{VAD} $=$ voice activity detector. Speech is active only about \mk{40\,\%} of
        the time, so \textbf{DTX} switches the transmitter off in the gaps.
  \item \mk{DTX saves handset battery and cuts interference to other cells.} Comfort noise
        is inserted at the far end so the line does not sound dead.
\end{itemize}

\lead{Channel coding \m{3}}
\begin{itemize}
  \item The 260 bits are \mk{sorted by how much they matter to speech quality}:
  \begin{itemize}
    \item \textbf{50 type Ia} bits, most important $\rightarrow$ \textbf{3 CRC parity} bits
          added.
    \item \textbf{132 type Ib} bits, next.
    \item \textbf{78 type II} bits, least important $\rightarrow$ \mk{no protection at all}.
  \end{itemize}
  \item Types Ia and Ib $+$ parity $+$ \textbf{4 tail zeros} $= 189$ bits go through a
        \mk{rate 1/2, constraint length $K=5$} convolutional encoder $\rightarrow$ 378 bits.
  \item $378 + 78 = \mathbf{456}$ bits per 20 ms $= \mathbf{22.8}$ kbps.
  \item \mk{Unequal error protection}: spending all the redundancy on the bits that would
        wreck the speech is what makes 13 kbps survive a fading channel.
\end{itemize}}

\vspace{2pt}
\sbsr{0.40}{%
\figT{c8_channel_coding.png}
\vspace{1pt}
{\color{sub}\footnotesize \textbf{260 bits in, 456 bits out.} \yr{SST Ch 8}}}{%
\lead{Interleaving, and why 8 slots}
\begin{itemize}
  \item A deep fade wipes out a \textbf{burst} of adjacent bits, which no convolutional code
        handles well.
  \item Spreading one speech frame over \mk{8 bursts in 8 different frames} turns one long
        error burst into eight short ones, \mk{one per code frame}, which the decoder can
        fix.
  \item Cost: \mk{about 40 ms of extra delay}, which is why the specification table lists
        interleaving depth as a delay figure.
  \item \mk{This is why a GSM call survives a moving car.} The channel decorrelates between
        bursts, so eight bursts are eight independent looks at the same speech frame, i.e.
        \textbf{time diversity} (Ch 5).
\end{itemize}}

\T{8.8 CDMA standards: IS-95 frequency and channel specification}

\Q{Frequency and channel specification of the IS-95 CDMA standard}{\m{2}\m{4}}
{\color{sub}\footnotesize \mk{no paper asks this alone}, but every IS-95 channel question
opens with it \gap$\cdot$\gap the CDMA \emph{principle}, near-far effect, encoding and
decoding are all \textbf{Ch 7}; this chapter is only the \textbf{standard}}

\sbsr{0.50}{%
\lead{IS-95, also called cdmaOne}
\begin{itemize}
  \item A \textbf{US standard} built on Qualcomm's technology, the first commercial CDMA
        cellular system.
  \item \textbf{FDD}, using two \mk{1.25 MHz simplex channels 45 MHz apart}.
  \begin{itemize}
    \item Reverse (uplink) \textbf{824--849 MHz}.
    \item Forward (downlink) \textbf{869--894 MHz}.
  \end{itemize}
  \item Chip rate \mk{1.2288 Mcps}, i.e. $128 \times 9.6$ kbps.
  \item \mk{Users in a cell and users in adjacent cells share the same radio channel}, so
        the frequency reuse factor is $N = 1$ and \textbf{no frequency planning is needed}.
        This is the single biggest planning advantage over GSM.
  \item $N=1$ is also what makes \textbf{soft handoff} possible: the mobile can be on the
        same carrier with two base stations at once (Ch 2).
  \item Variable-rate vocoder: \mk{9.6 / 4.8 / 2.4 / 1.2 kbps}. Later \textbf{EVRC} gives
        13 kbps quality at the 8 kbps rate.
\end{itemize}}{%
\lead{The three code families \m{3}}
\begin{itemize}
  \item \textbf{Walsh codes}, 64 rows of a $64\times64$ Hadamard matrix.
  \begin{itemize}
    \item \emph{Forward link}: one Walsh code per \textbf{logical channel}, so they separate
          \mk{channels}.
    \item \emph{Reverse link}: 64-ary orthogonal modulation, so they separate
          \mk{groups of six symbols}, not users. \mk{Learn this asymmetry, it is examinable.}
  \end{itemize}
  \item \textbf{Short PN codes}, a quadrature pair of period $2^{15}$ chips.
  \begin{itemize}
    \item Applied on the I and Q arms after Walsh spreading.
    \item Each cell or sector uses one of \mk{512 phase offsets}, so the short code
          \textbf{identifies the sector}.
  \end{itemize}
  \item \textbf{Long PN code}, period $2^{42}-1$, generated from a 42-bit mask.
  \begin{itemize}
    \item \emph{Forward}: scrambles and encrypts the traffic, mask derived from the user's
          \textbf{ESN}.
    \item \emph{Reverse}: \mk{it is the channelising code} $-$ each mobile is separated by
          its own long-code phase.
  \end{itemize}
\end{itemize}

\lead{Power control, because CDMA lives or dies by it \m{2}}
\begin{itemize}
  \item \textbf{Open loop}: the mobile measures the received \textbf{pilot} strength and
        sets its own transmit power inversely.
  \item \textbf{Closed loop}: the base station measures each mobile and sends
        \textbf{power control bits} at \mk{800 Hz}, $\pm 1$ dB per bit.
  \item Without it the \textbf{near-far problem} (Ch 7 \S7.6) destroys capacity.
\end{itemize}}

\T{8.9 Forward CDMA (IS-95) channel \; \small\mk{the standard 4 to 6 mark drawing}}

\Q{Draw and explain the forward CDMA (IS-95) channel; what is a channelization code?}{\m{4}\m{6}}
{\color{sub}\footnotesize \tS{6} \yr{\textbf{\texttt{80 Bh}} \m{4}, \textbf{80 Ch} \m{6},
\textbf{78 Ch} \m{4}, 73 Ma \m{4} (channelization code $+$ forward channels),
\textbf{73 Bh} \m{6} (pilot and sync channels only), \texttt{81 Ba} \m{4} (reverse, \S8.10)}
\gap$\cdot$\gap \mk{73 Bh wants only the pilot and sync blocks}, drawn out in full}

\sbsr{0.44}{%
\figT{c8_walsh_gen.png}
\vspace{1pt}
{\color{sub}\footnotesize \textbf{Channelization code $=$ Walsh code} \m{2}, built by the
Hadamard recursion. Row $n$ of the matrix is $W_n$. Worked as a numerical in \textbf{Ch 7
companion \S4}. \yr{SST Ch 8}}

\vspace{3pt}
\lead{Why Walsh codes work here}
\begin{itemize}
  \item All 64 rows are \textbf{mutually orthogonal}: map $0 \rightarrow -1$,
        $1 \rightarrow +1$ and the term-by-term product of any two rows sums to \textbf{0}.
  \item So a receiver that correlates with $W_n$ recovers channel $n$ and sees \mk{exactly
        zero} from the other 63, provided they arrive \textbf{in step}.
  \item They stay in step on the forward link because \mk{one base station transmits them
        all}. That is why the reverse link cannot use them for channelisation.
\end{itemize}}{%
\lead{The four forward logical channels \m{4}}
\begin{enumerate}
  \item \textbf{Pilot channel, $W_0$.}
  \begin{itemize}
    \item Carries \mk{no baseband information at all}: an all-zero stream, spread by $W_0$
          which is itself all zeros, so what goes out \textbf{is the short PN pair itself}
          at the sector's offset.
    \item Transmitted \textbf{continuously}. Gives the mobile \textbf{timing acquisition and
          a phase reference} for coherent detection.
    \item Its $E_c/I_0$ is what the mobile measures to decide \mk{which sector is strongest},
          i.e. it drives handoff.
  \end{itemize}
  \item \textbf{Sync channel, $W_{32}$.}
  \begin{itemize}
    \item \mk{1.2 kbps} of real baseband data, convolutionally coded and interleaved.
    \item Carries the sync channel message: \textbf{system time, long-code state, the
          sector's PN offset, and the paging channel rate} (PRAT).
    \item Uses the \textbf{same PN offset as the pilot}, so once the pilot is acquired the
          sync channel is readable at once.
  \end{itemize}
\end{enumerate}}

\vspace{2pt}
\sbsr{0.49}{%
\begin{enumerate}
  \setcounter{enumi}{2}
  \item \textbf{Paging channels, $W_1$ to $W_7$}, up to seven.
  \begin{itemize}
    \item \mk{4.8 or 9.6 kbps}, the rate announced by the sync message.
    \item Overhead messages, \textbf{pages}, and \textbf{traffic channel assignment} during
          call setup.
    \item A mobile monitors \textbf{one} paging channel, and ignores it once a traffic
          channel is up.
  \end{itemize}
\end{enumerate}}{%
\begin{enumerate}
  \setcounter{enumi}{3}
  \item \textbf{Forward traffic channels}, the remaining Walsh codes.
  \begin{itemize}
    \item User voice and data at \mk{1.2 / 2.4 / 4.8 / 9.6 kbps}.
    \item Structure is the paging channel plus one addition: \mk{multiplexed power control
          bits at 800 bps}, punctured into the stream.
    \item Scrambled by the long code with that user's \textbf{ESN}-derived mask.
  \end{itemize}
\end{enumerate}}

\vspace{2pt}
\sbsr{0.63}{%
\figT{c8_fwd_cdma.png}
\vspace{1pt}
{\color{sub}\footnotesize \textbf{Forward traffic channel, the block diagram to reproduce.}
Rate 1/2 $K=9$ convolutional encoder $\rightarrow$ symbol repetition to 19.2 ksps
$\rightarrow$ block interleaver $\rightarrow$ long-code scrambling $\rightarrow$ MUX with the
800 bps power control bits $\rightarrow$ Walsh cover $W_i$ $\rightarrow$ I/Q short PN
$\rightarrow$ QPSK. \yr{SST Ch 8}}}{%
\lead{The modulator, and how the channels combine \m{2}}
\begin{itemize}
  \item Every logical channel first passes a \textbf{gain control} block, which sets how much
        of the sector's power that channel gets.
  \begin{itemize}
    \item Traffic channel gains \mk{change continuously}, driven by forward power control.
    \item Pilot gain is fixed and deliberately generous, typically 15 to 20\,\% of total
          power, because everything else depends on it. \added
  \end{itemize}
  \item After gain setting all channels are \textbf{coherently added} into one composite
        spread-spectrum signal.
  \item The sum is split into \textbf{I and Q}, each spread by its short PN sequence,
        up-converted on quadrature carriers and added $\rightarrow$ one \textbf{QPSK}
        passband signal.
  \item \mk{One 1.25 MHz carrier therefore carries all 64 code channels of the sector at
        once.}
\end{itemize}}

\vspace{2pt}
\sbs{%
\figT{c8_pilot_ch.png}
\vspace{1pt}
{\color{sub}\footnotesize \textbf{Pilot channel} \m{3}. All zeros in, PN out. No encoder, no
interleaver, nothing to decode. \yr{SSD, IS-95 channels}}}{%
\figT{c8_sync_ch.png}
\vspace{1pt}
{\color{sub}\footnotesize \textbf{Sync channel} \m{3}. 1.2 kbps $\rightarrow$ rate 1/2
$\rightarrow$ 2.4 ksps $\rightarrow$ repetition to 4.8 ksps $\rightarrow$ interleave
$\rightarrow$ $W_{32}$ $\rightarrow$ I/Q PN. \yr{SSD, IS-95 channels}}}

\vspace{2pt}
\sbs{%
\figT{c8_paging_ch.png}
\vspace{1pt}
{\color{sub}\footnotesize \textbf{Paging channel}. Same shape as the sync channel plus
long-code scrambling through a decimator. \yr{SSD, IS-95 channels}}}{%
\figT{c8_fwd_traffic_ch.png}
\vspace{1pt}
{\color{sub}\footnotesize \textbf{Forward traffic channel}. The paging channel plus the
\mk{MUX that punctures in the 800 bps power control bits}. \yr{SSD, IS-95 channels}}}

\T{8.10 Reverse CDMA (IS-95) channel}

\Q{Draw the reverse CDMA (IS-95) channel}{\m{4}}
{\color{sub}\footnotesize \yr{\texttt{81 Ba} \m{4}} \gap$\cdot$\gap
\mk{counted inside the \tS{6} forward/reverse total of \S8.9} \gap$\cdot$\gap
the marks are usually won on \mk{why the reverse link is built differently}, not on the
drawing}

\sbsr{0.50}{%
\lead{Two channels only, and no pilot}
\begin{itemize}
  \item The reverse link supports \textbf{access channels} and \textbf{traffic channels}.
        That is all.
  \item \mk{There is no reverse pilot} in IS-95: it is impractical for every mobile to
        broadcast its own pilot sequence.
  \item No pilot $\rightarrow$ no phase reference $\rightarrow$ the reverse link is
        \textbf{non-coherent}.
  \item Non-coherent $\rightarrow$ \mk{Walsh functions cannot channelise}, because mobiles
        cannot be held chip-aligned with each other.
  \item Instead \textbf{each mobile is identified by its own long PN sequence}, set by a
        mask derived from its \textbf{ESN}.
  \item Walsh functions are still used, but for \mk{64-ary orthogonal modulation}: every six
        code symbols are sent as \emph{one} of the 64 Walsh functions, and the base station
        decides which one by correlating against all 64.
  \item \mk{Forward: Walsh separates channels. Reverse: Walsh represents symbols.}
\end{itemize}}{%
\figT{c8_rev_cdma.png}
\vspace{1pt}
{\color{sub}\footnotesize \textbf{Reverse traffic and access channel generation.} Rate 1/3
$K=9$ encoder $\rightarrow$ repetition to 28.8 ksps $\rightarrow$ interleaver
$\rightarrow$ 64-ary orthogonal modulator, 307.2 kcps $\rightarrow$ data burst randomizer
$\rightarrow$ long PN $\rightarrow$ I/Q short PN with a \mk{half-chip delay on Q}, giving
OQPSK. \yr{SST Ch 8}}}

\vspace{2pt}
\sbs{%
\lead{1. Access channel}
\begin{itemize}
  \item Used when the mobile \textbf{has no traffic channel assigned}: to register,
        originate calls, respond to pages and orders, and send overhead messages.
  \item Baseband rate is \textbf{fixed at 4.8 kbps}.
  \item \mk{Rate 1/3 convolutional encoder}, not 1/2 as on the forward link. The reverse is
        the weaker of the two links (a handset transmits under a watt), so it is given more
        redundancy.
  \item Symbol repetition doubles it to \mk{28.8 ksps}, then block interleaving, then the
        64-ary orthogonal modulator.
\end{itemize}}{%
\lead{2. Reverse traffic channel}
\begin{itemize}
  \item Carries user voice, data and signalling at \mk{1.2 / 2.4 / 4.8 / 9.6 kbps}.
  \item Structure is the access channel plus one block: the \textbf{data burst randomizer}.
  \item \mk{The randomizer is the whole point.} When the vocoder drops to a lower rate the
        repeated symbols are redundant, so the randomizer \textbf{gates them off} instead of
        transmitting them.
  \item Result: the handset transmits only in pseudo-random bursts $\rightarrow$ lower
        average power, longer battery life and \mk{less interference into every other user},
        which is capacity.
  \item The forward link exploits the same voice activity differently: it \textbf{lowers the
        energy per symbol} rather than gating off, since the base station must keep the
        composite waveform continuous.
\end{itemize}}

\vspace{2pt}
\sbs{%
\figT{c8_access_ch.png}
\vspace{1pt}
{\color{sub}\footnotesize \textbf{Access channel} \m{2}. Note the fixed 4.8 kbps input and
the absence of a data burst randomizer. \yr{SSD, IS-95 channels}}}{%
\figT{c8_rev_traffic_ch.png}
\vspace{1pt}
{\color{sub}\footnotesize \textbf{Reverse traffic channel} \m{2}. The one added block is the
\mk{data burst randomizer}. \yr{SSD, IS-95 channels}}}

\T{8.11 GSM against CDMA, and CDMA against LTE}

\Q{Compare GSM and CDMA standards}{\m{3}\m{4}}
{\color{sub}\footnotesize \yr{\textbf{76 Bh} \m{3}, 71 Ma \m{4}} \gap$\cdot$\gap
71 Ma phrases it as \mk{``what is GSM and CDMA standard''} before asking for the GSM
architecture, so the comparison is the short half}

\par\vspace{2pt}
\begin{tabularx}{\textwidth}{@{}L{3.0cm}XX@{}}
\toprule
\hrow \thead{Point} & \thead{GSM} & \thead{CDMA (IS-95)} \\
\midrule
Multiple access & TDMA over FDMA & CDMA, direct sequence \\
Duplexing & FDD & FDD \\
Carrier bandwidth & 200 kHz & 1.25 MHz \\
Users per carrier & 8 (16 at half rate) & about 55 to 61 code channels \\
Frequency reuse & $N = 3, 4, 7, 9$; \mk{needs planning} & \mk{$N = 1$, no planning} \\
Modulation & 0.3 GMSK & QPSK forward, OQPSK reverse \\
Speech coder & RPE-LTP 13 kbps, fixed & variable rate 9.6 down to 1.2 kbps \\
Handoff & \mk{hard} (break before make) & \mk{soft and softer} (make before break) \\
Power control & slow, on SACCH & \mk{fast, 800 Hz closed loop}, essential \\
Equalizer & \mk{required}, 26-bit training sequence & not needed; \mk{RAKE receiver} instead \\
Capacity limit & hard: run out of slots & \mk{soft: interference limited}, degrades gently \\
Subscriber identity & \mk{SIM card}, moves between handsets & tied to the handset ESN \\
Security & A3/A5/A8 over the air & \mk{inherent} in the long PN code, plus ciphering \\
\bottomrule
\end{tabularx}

\Q{Compare the system architecture of CDMA with LTE; mention the function of the entities}{\m{3+5}}
{\color{sub}\footnotesize \tP{1} \yr{74 Ma \m{3+5}} \gap$\cdot$\gap \added \;
\mk{no lecture deck covers LTE at all} \gap$\cdot$\gap 8 marks, so it is worth learning the
seven LTE boxes by name}

\sbs{%
\lead{CDMA (IS-95 / cdma2000) architecture}
\begin{itemize}
  \item \textbf{MS} $\rightarrow$ \textbf{BTS} $\rightarrow$ \textbf{BSC} $\rightarrow$
        \textbf{MSC} for voice, or \textbf{PDSN} for packet data.
  \item \textbf{BTS}: radio, Walsh covering, power control commands.
  \item \textbf{BSC}: radio resource control, \mk{soft handoff combining}, and the vocoders.
  \item \textbf{MSC}: circuit switching to the PSTN.
  \item \textbf{PCF} $+$ \textbf{PDSN}: the packet data path, bolted on later.
  \item \textbf{HLR / VLR / AuC}: same roles as in GSM.
  \item \mk{Hierarchical and circuit-first}: voice and data ride separate cores.
\end{itemize}}{%
\lead{LTE (EPS) architecture}
\begin{itemize}
  \item \textbf{UE} $\rightarrow$ \textbf{eNodeB} $\rightarrow$ \textbf{EPC}
        (\textbf{MME}, \textbf{S-GW}, \textbf{P-GW}) $\rightarrow$ external IP networks.
  \item \textbf{eNodeB}: radio plus \mk{all the radio resource control that a BSC used to
        do}. There is no controller node. eNodeBs talk to each other directly over
        \textbf{X2}.
  \item \textbf{MME}, control plane only: attach, authentication, bearer setup, mobility and
        paging.
  \item \textbf{S-GW}: user-plane anchor, forwards packets during inter-eNodeB handover.
  \item \textbf{P-GW}: IP anchor to the outside world, allocates the IP address, enforces
        policy, does charging.
  \item \textbf{HSS}: subscriber database, \mk{HLR $+$ AuC in one}.
  \item \textbf{PCRF}: policy and charging rules.
  \item \mk{Flat and all-IP}: two nodes from handset to internet, and no circuit domain at
        all.
\end{itemize}}

\vspace{2pt}
\lead{The three differences to write down \m{3}}
\begin{itemize}
  \item \textbf{Topology}: CDMA is BTS--BSC--MSC, three tiers. LTE is eNodeB--EPC, two tiers.
        \mk{The BSC's job moved into the base station}, which cuts handover latency.
  \item \textbf{Bearer}: CDMA is circuit-switched voice with packet data added later. LTE is
        \mk{all-IP, packet only}; voice comes back as VoLTE over IMS.
  \item \textbf{Air interface}: 1.25 MHz DS-CDMA against \mk{OFDMA down, SC-FDMA up, with
        MIMO}, on a scalable 1.4 to 20 MHz carrier. Peak rate rises from about 3 Mbps to
        100 Mbps and beyond.
\end{itemize}

\T{8.12 Evolution of wireless telephone systems \; \pillo{gd}{syllabus, no PYQ yet}}

\Q{AMPS, CT2, DECT, PHS, PACS, IS-54 / IS-136, IS-94, IS-95: the standards named in the syllabus}{}
{\color{sub}\footnotesize \pill{gdL}{gd}{syllabus, no PYQ yet} \gap$\cdot$\gap
syllabus item \textbf{8(a)}, \mk{never examined in 22 papers} \gap$\cdot$\gap
the \emph{generations} question (1G to 5G) is \textbf{Ch 1 \S1.3} and is asked every year;
this band is the \emph{systems} behind those generations \gap$\cdot$\gap
\mk{read the three tables, do not memorise them}}

\sbs{%
\lead{First generation, analog cellular}
\begin{itemize}
  \item All \textbf{FM voice with FSK signalling, FDMA/FDD}, one carrier per user for the
        whole call.
  \item \textbf{AMPS}, the Americas: 824--849 / 869--894 MHz, 30 kHz channels,
        \mk{832 channels}.
  \item \textbf{TACS / ETACS}, UK and Europe: 25 kHz channels, up to 1240 channels.
  \item \textbf{NMT-450 / NMT-900}, the Nordic countries, 25 and 12.5 kHz.
  \item \textbf{C-450} Germany, \textbf{RTMS} Italy, \textbf{Radiocom 2000} France,
        \textbf{NTT} and \textbf{JTACS/NTACS} Japan.
  \item \mk{Every country built its own}, which is exactly the fragmentation GSM was created
        to end.
\end{itemize}}{%
\figT{c8_analog_systems.png}
\vspace{1pt}
{\color{sub}\footnotesize \yr{Rappaport Table 10.10}}}

\vspace{2pt}
\sbs{%
\figT{c8_cordless_std.png}
\vspace{1pt}
{\color{sub}\footnotesize \textbf{Digital cordless standards.} \yr{Rappaport Table 10.11}}}{%
\figT{c8_2g_std.png}
\vspace{1pt}
{\color{sub}\footnotesize \textbf{Second generation digital cellular.} \yr{Rappaport
Table 10.12}}}

\vspace{2pt}
\sbs{%
\lead{Cordless and PCS, the short-range family}
\begin{itemize}
  \item \textbf{CT2}, Europe and Canada, 864--868 MHz, \mk{TDD}, 100 kHz carriers,
        \textbf{one channel per carrier}, 72 kbps GFSK, ADPCM speech.
  \item \textbf{DECT}, Europe, 1880--1900 MHz, TDD, \mk{1728 kHz carriers with 12 channels
        each}, 1152 kbps GFSK. The one that survived, as the office and home cordless phone.
  \item \textbf{PHS}, Japan, 1895--1918 MHz, TDD, 300 kHz, 4 channels per carrier,
        384 kbps $\pi/4$-DQPSK.
  \item \textbf{PACS}, USA, \mk{FDD or TDD}, 300 kHz, 8 or 4 channels per carrier,
        384 kbps $\pi/4$-QPSK. Designed for licensed PCS microcells.
  \item \textbf{IS-94} is the odd one out: it is not cordless but the standard for
        \mk{AMPS-based private and in-building microcell systems}, i.e. cellular hardware
        reused as a premises system. \added
  \item \mk{All of them are TDD, low power and short range}, which is what separates a
        cordless standard from a cellular one.
\end{itemize}}{%
\lead{Second generation, digital cellular}
\begin{itemize}
  \item \textbf{GSM} (1990): TDMA/FDMA/FDD, 200 kHz, GMSK $BT = 0.3$, 270.833 kbps,
        RPE-LTP at 13 kbps, \mk{1.35 bps/Hz}.
  \item \textbf{IS-54}, later \textbf{IS-136}, also called \textbf{USDC} or D-AMPS (1991):
        \mk{30 kHz channels, deliberately the same as AMPS} so the two could share a band,
        3 users per carrier, $\pi/4$-DQPSK at 48.6 kbps, VSELP 7.95 kbps.
  \item \textbf{PDC}, Japan (1993): 25 kHz, $\pi/4$-DQPSK at 42 kbps, VSELP 6.7 kbps,
        \mk{1.68 bps/Hz}, the most spectrally efficient of the three.
  \item \textbf{IS-95} (1993): the CDMA outlier, 1.25 MHz carriers, $N=1$. \S8.8.
  \item \mk{Three of the four are TDMA.} CDMA was the minority choice in 2G and the majority
        choice in 3G.
\end{itemize}}

\T{8.13 WiFi, WiMAX, UMTS, EVDO, UMB, LTE and recent trends \; \pillo{gd}{no lecture note covers this}}

\Q{Write short notes on WiMAX / LTE}{\m{3}\m{4}\m{5}}
{\color{sub}\footnotesize \tF{3} \yr{\textbf{78 Ch} \m{4} (WiMAX), 73 Ma \m{5} (WiMAX),
\textbf{72 Ash} \m{3} (WiMAX) $+$ \m{3} (LTE)} \gap$\cdot$\gap \added \;
\mk{absent from all three lecture decks} \gap$\cdot$\gap
\textbf{WiFi and WLAN are answered in Ch 7 \S7.11}, not here}

\sbs{%
\lead{WiMAX \m{4}}
\begin{itemize}
  \item \textbf{WiMAX} $=$ worldwide interoperability for microwave access, the trade name
        for \textbf{IEEE 802.16}. A \mk{WMAN}, one step above WiFi's WLAN.
  \item Two versions worth naming:
  \begin{itemize}
    \item \textbf{802.16d} (2004), \textbf{fixed} WiMAX, OFDM.
    \item \textbf{802.16e} (2005), \textbf{mobile} WiMAX, \mk{scalable OFDMA}, supports
          handover at vehicular speed.
  \end{itemize}
  \item Bands: \mk{2--11 GHz for NLOS}, 10--66 GHz for LOS backhaul. Scalable channel
        bandwidth 1.25 to 20 MHz.
  \item Throughput about \mk{75 Mbps fixed}, 15 Mbps mobile; range 5--8 km NLOS, up to
        50 km LOS.
  \item \mk{The MAC is the real difference from WiFi}: connection-oriented and
        \textbf{scheduled by the base station}, not contended. Every subscriber gets a
        grant, so QoS classes (UGS, rtPS, nrtPS, best effort) are enforceable.
  \item Topologies: point-to-multipoint, and mesh.
  \item \textbf{Uses}: last-mile broadband where cable is uneconomic, backhaul for WiFi
        hotspots and cell sites, rural and disaster-recovery connectivity.
  \item \textbf{(-)}: needs licensed spectrum, expensive base stations, and it
        \mk{lost the 4G market to LTE} because operators already owned 3GPP cores.
\end{itemize}}{%
\lead{LTE \m{3}}
\begin{itemize}
  \item \textbf{LTE} $=$ long term evolution, 3GPP Release 8 (2008). The 4G standard that
        won.
  \item \textbf{OFDMA} on the downlink, \textbf{SC-FDMA} on the uplink.
        \mk{SC-FDMA is chosen purely for its low PAPR}, which is what a battery-powered
        handset amplifier needs (Ch 4 \S4.6).
  \item \textbf{MIMO} up to $4\times4$, scalable bandwidth
        \mk{1.4 / 3 / 5 / 10 / 15 / 20 MHz}, both FDD and TDD.
  \item Release 8 peaks at about \mk{300 Mbps down, 75 Mbps up}.
  \item \textbf{Flat all-IP} architecture, \S8.11. \mk{No circuit-switched domain}: voice is
        VoLTE over IMS.
  \item \textbf{LTE-Advanced} (Rel 10) adds \textbf{carrier aggregation}, $8\times8$ MIMO,
        relay nodes and CoMP, reaching \mk{1 Gbps}.
\end{itemize}

\lead{The other names on the syllabus \m{2}}
\begin{itemize}
  \item \textbf{UMTS}: 3GPP's 3G. \mk{W-CDMA, 5 MHz carriers, 3.84 Mcps}. Radio network is
        NodeB $+$ RNC, reusing the GSM/GPRS core. HSPA later takes it to 14.4 Mbps.
  \item \textbf{CDMA-EVDO}: cdma2000 \mk{1xEV-DO, ``evolution, data optimised''}. A
        \emph{separate} 1.25 MHz data-only carrier with a time-multiplexed forward link and
        rate adaptation. Rev 0 gives 2.4 Mbps, Rev A 3.1 Mbps down / 1.8 Mbps up.
  \item \textbf{UMB}: ultra mobile broadband, 3GPP2's OFDMA answer to LTE.
        \mk{Abandoned in 2008} when Qualcomm stopped development and the CDMA operators moved
        to LTE. Worth one honest line: it is on the syllabus but was never deployed. \added
\end{itemize}}

\vspace{2pt}
\sbsr{0.55}{%
\lead{Recent trends \m{2}}
\begin{itemize}
  \item \textbf{5G NR}: sub-6 GHz plus \mk{mmWave}, massive MIMO and beamforming, network
        slicing, three service classes \textbf{eMBB} / \textbf{URLLC} / \textbf{mMTC}.
  \item \textbf{IoT}: NB-IoT and LTE-M, narrowband, low power, huge device counts.
  \item \textbf{D2D} device-to-device, \textbf{C-RAN} centralised baseband, edge computing.
  \item \mk{Cross-refer Ch 1 \S1.4}, which answers the ``trends beyond 3G'' question in
        full. Do not repeat it here.
\end{itemize}}{%
\par\vspace{2pt}
\begin{tabularx}{\linewidth}{@{}L{1.7cm}XXX@{}}
\toprule
\hrow \thead{} & \thead{WiFi} & \thead{WiMAX} & \thead{LTE} \\
\midrule
Standard & 802.11 & 802.16 & 3GPP Rel 8 \\
Scale & WLAN, 100 m & WMAN, km & cellular WAN \\
Spectrum & \mk{unlicensed} & licensed & licensed \\
MAC & \mk{contended}, CSMA/CA & scheduled & scheduled \\
QoS & weak & \mk{built in} & built in \\
Mobility & nomadic & vehicular (802.16e) & \mk{full handover} \\
\bottomrule
\end{tabularx}
\vspace{2pt}
{\color{sub}\footnotesize \mk{If a question asks you to ``compare'', this table is the
answer.} The row that actually decides all the others is \textbf{contended against
scheduled}. \added}}

\T{8.14 Regulatory issues \; \small\mk{the highest-frequency topic in this chapter}}

\Q{Describe regulatory issues in wireless communication: spectrum allocation, pricing, licensing, tariffs, interconnection}{\m{4}\m{5}}
{\color{sub}\footnotesize \tS{9} \yr{\textbf{\texttt{81 Bh}} \m{4} (licensing $+$ allocation),
\texttt{80 Ba} \m{4}, \textbf{\texttt{79 Bh}} \m{4} (spectrum regulations),
\textbf{76 Bh} \m{5}, \textbf{74 Bh} \m{5}, \textbf{73 Bh} \m{5} (spectral licensing),
\textbf{80 Ch} \m{2+4}, \textbf{79 Ch} \m{2+3}, 70 Ma \m{4}} \gap$\cdot$\gap
\mk{80 Ch and 79 Ch ask ``significance of spectrum management'' first} \m{2}, then the
\textbf{functions of the regulator}}

\sbsr{0.50}{%
\lead{Why regulation exists at all \m{2}}
\begin{itemize}
  \item Radio spectrum is a \textbf{finite, non-storable, shared} natural resource.
  \item Interference is \mk{non-excludable}: two users on the same frequency, in the same
        place, at the same time both fail. Neither can protect itself.
  \item So a government must decide \textbf{who may use which band, where, and on what
        terms}. That is spectrum management.
  \item \textbf{Significance of spectrum management} \m{2}:
  \begin{itemize}
    \item Interference-free coexistence of services.
    \item \mk{Efficient use} of a scarce resource, and no hoarding.
    \item \textbf{International harmonisation} $\rightarrow$ roaming, economies of scale in
          handsets.
    \item Protected bands for \textbf{aviation, defence, emergency and scientific} use.
    \item Predictable rules $\rightarrow$ operators will invest.
  \end{itemize}
\end{itemize}

\lead{1. Spectrum allocation}
\begin{itemize}
  \item \textbf{ITU-R} divides the world into \mk{three Regions} and allocates bands to
        \emph{services} at the World Radiocommunication Conference.
  \item The \textbf{national regulator} then assigns those bands to \emph{operators}:
        \textbf{NTA} in Nepal, FCC in the USA, CEPT/ETSI in Europe.
  \item Assignment methods: \textbf{administrative} (a ``beauty contest'' on merit),
        \textbf{auction}, lottery, first-come first-served, or left
        \mk{unlicensed as a commons} (the ISM bands WiFi and Bluetooth live in).
\end{itemize}}

\vspace{2pt}
\sbs{%
\lead{2. Spectrum pricing}
\begin{itemize}
  \item Charges an operator for holding spectrum: an \textbf{annual royalty} per MHz per head
        of population, or the one-off \textbf{auction price}.
  \item Aim: \mk{make idle spectrum expensive}, so it moves to whoever will use it best.
  \item \textbf{Risk}: overpricing is passed straight to the subscriber as tariff, and slows
        rural rollout. \mk{The regulator has to balance revenue against coverage.}
\end{itemize}}{%
\lead{3. Licensing}
\begin{itemize}
  \item Decides \textbf{who may operate}, over what area, for how long, and on renewal terms.
  \item Attaches obligations: \textbf{coverage targets, quality of service, rural
        rollout}, lawful interception.
  \item \textbf{Spectrum caps} stop one operator cornering the band and keep the market
        competitive.
  \item \textbf{Type approval}: only equipment that meets emission limits may be connected.
  \item Technology-neutral licences let an operator refarm a 2G band for LTE without a new
        licence.
\end{itemize}}

\vspace{2pt}
\sbs{%
\lead{4. Tariff regulation}
\begin{itemize}
  \item Regulator \textbf{approves tariffs}, and classifies them as \mk{fixed, maximum, or
        non-regulated}.
  \item Price caps and floors protect subscribers from a dominant operator and rivals from
        \textbf{predatory pricing}.
  \item Tariffs are usually deregulated as competition arrives.
\end{itemize}}{%
\lead{5. Interconnection}
\begin{itemize}
  \item Every licensee must be able to \textbf{terminate calls on every other network},
        otherwise the first operator to scale wins by default.
  \item Regulator sets \mk{cost-oriented interconnection charges}, the points of
        interconnection, and \textbf{number portability}.
  \item It also \textbf{settles interconnection disputes}, which is why dispute resolution
        appears in the list of functions below.
\end{itemize}}

\vspace{2pt}
\sbsr{0.52}{%
\lead{Functions of the regulator \m{4} \; \mk{the list 80 Ch and 79 Ch want}}
\begin{enumerate}
  \item \textbf{Inspection and investigation} of the affairs of telecom operators and
        internet service providers.
  \item \textbf{Settle disputes} between licensees, or between a licensee and a customer,
        relating to telecommunications service.
  \item \textbf{Determine the quality and standard} of the machinery, equipment and
        facilities used, and of the service itself.
  \item \textbf{Licensing} of telecom operators and internet service providers.
  \item \textbf{Approval of tariffs}: fixed tariffs, maximum tariffs and non-regulated
        tariffs.
\end{enumerate}
\vspace{2pt}
{\color{sub}\footnotesize \yr{SST Ch 8, the NTA regulatory approach} \gap$\cdot$\gap
\mk{write these five and the 4 marks are secure.} Add spectrum allocation and monitoring as
a sixth if the question says ``spectrum management''.}}{%
\lead{Nepal, if a local example is wanted \added}
\begin{itemize}
  \item \textbf{NTA} $=$ Nepal Telecommunications Authority, constituted under the
        \textbf{Telecommunications Act 2053}.
  \item Licence categories it issues include \textbf{basic telephony, cellular mobile,
        internet service, VSAT and rural VSAT, network service providers, international
        trunk telephone} and rural ISP.
  \item \mk{Naming NTA and the Act is enough} for a Nepali paper. Do not attempt licence
        counts, they change every year.
\end{itemize}

\lead{Answer skeleton \m{5}}
\begin{enumerate}
  \item One line: spectrum is finite and shared, so it must be regulated. \m{1}
  \item Name the \textbf{five heads}: allocation, pricing, licensing, tariff, interconnection.
        \m{1}
  \item Two sentences on each head. \m{2.5}
  \item Close with ITU-R above and the national regulator (NTA) below. \m{0.5}
\end{enumerate}}

\T{8.15 Last-minute recall for this chapter}

\sbsr{0.52}{%
\lead{Must memorise}
\begin{itemize}
  \item GSM 900: \textbf{890--915} up, \textbf{935--960} down, \mk{45 MHz} apart, 200 kHz
        ARFCN, \mk{124} duplex channels, 8 users each.
  \item \textbf{0.3 GMSK}, \mk{270.833 kbps} per carrier, 33.854 kbps per user, 24.7 kbps
        max user data.
  \item Slot \mk{156.25 bits, 576.9 $\mu$s}; frame \mk{8 slots, 4.615 ms}; bit
        \mk{3.692 $\mu$s}.
  \item Normal burst $3 + 57 + 1 + 26 + 1 + 57 + 3 + 8.25$.
  \item Multiframe \mk{26 (traffic, 120 ms)} or \mk{51 (control, 235 ms)};
        superframe \mk{1326 frames, 6.12 s}; hyperframe \mk{2048 superframes, 3 h 28 min}.
  \item Frame efficiency \mk{74.24\,\%} $\rightarrow$ Ch 7 companion \S3.
  \item Speech \mk{260 bits / 20 ms $=$ 13 kbps} $\rightarrow$ coded to
        \mk{456 bits $=$ 22.8 kbps} $\rightarrow$ interleaved over \mk{8 slots}.
  \item Channel coding split \mk{50 Ia $+$ 132 Ib $+$ 78 II}, rate 1/2, $K=5$.
  \item Control channels: \mk{BCH $=$ BCCH, FCCH, SCH} $\cdot$
        \mk{CCCH $=$ PCH, RACH, AGCH} $\cdot$ \mk{DCCH $=$ SDCCH, SACCH, FACCH}.
        SACCH every \mk{13th} frame; FACCH \mk{steals} a TCH burst.
  \item IS-95: \mk{1.25 MHz}, 824--849 / 869--894, chip rate \mk{1.2288 Mcps},
        \mk{$N=1$}, power control at \mk{800 Hz}.
  \item Forward channels: \mk{pilot $W_0$, sync $W_{32}$, paging $W_1$--$W_7$, traffic}.
        Reverse: \mk{access and traffic only, no pilot}.
  \item Forward encoder \mk{rate 1/2 $K=9$}; reverse \mk{rate 1/3 $K=9$} $+$ 64-ary
        orthogonal modulation $+$ data burst randomizer.
\end{itemize}}{%
\lead{Most repeated, in order}
\begin{enumerate}
  \item Regulatory issues, or significance of spectrum management \tS{9}
  \item GSM traffic and control channels \tS{8}
  \item GSM system architecture, or the NSS alone \tS{7}
  \item Forward and reverse CDMA (IS-95) channels \tS{6}
  \item GSM frame structure and hierarchy \tF{5}
  \item WiMAX short note \tF{3}
  \item Signal processing in GSM \tP{2} \gap$\cdot$\gap
        how channels are used while making a call \tF{2}
  \item CDMA against LTE architecture \tP{1} \gap$\cdot$\gap
        GSM against CDMA, specifications of GSM, LTE short note
\end{enumerate}
\vspace{2pt}
{\color{sub}\footnotesize \mk{The best hour you can spend on this chapter} is the
\textbf{control channel tree}: nine channels, one line each, plus which direction each one
points. Eight of the twenty-two papers ask for it directly, and it also supplies most of the
GSM-call answer in \S8.6.}

\vspace{2pt}
{\color{sub}\footnotesize \mk{No numerical companion for this chapter.} No paper from 2070 to
2082 sets a calculation in Chapter 8. The one number the syllabus touches, \textbf{GSM frame
efficiency}, is examined as a Chapter 7 multiple-access question and is worked in full in
\textbf{Ch 7 $-$ Numerical Problems \& Solutions, \S3}, including the 74 Ma
``cannot reach 100\,\%'' proof.}

\vspace{2pt}
{\color{sub}\footnotesize \mk{Two traps in this chapter}: \textbf{BCH} is the broadcast
channel here and a block code in Ch 6; and \textbf{FCCH} (GSM frequency correction channel)
is not \textbf{FCC} (forward control channel, Ch 2). Read the question's context before you
expand either.}}
